% Pi Coding Agent vertical tutorial PDF header.
% Rendered with Pandoc + XeLaTeX.

\usepackage{fontspec}
\usepackage{xcolor}
\usepackage{geometry}
\usepackage{titlesec}
\usepackage{fancyhdr}
\usepackage{microtype}
\usepackage{enumitem}
\usepackage{booktabs}
\usepackage{longtable}
\usepackage{array}
\usepackage{framed}
\usepackage{hyperref}

% Keep the visual language aligned with the tutorial decks.
\definecolor{piaccent}{HTML}{D77600}
\definecolor{piblue}{HTML}{4A90D9}
\definecolor{pimut}{HTML}{7A7A78}
\definecolor{piblack}{HTML}{0A0A0A}
\definecolor{pibody}{HTML}{242424}
\definecolor{picode}{HTML}{F3F0EB}
\definecolor{piline}{HTML}{DDD7CE}

\geometry{
  paperwidth=8.27in,
  paperheight=11.69in,
  top=0.78in,
  bottom=0.72in,
  inner=0.86in,
  outer=0.72in,
  headheight=15pt,
  headsep=0.18in,
  footskip=0.36in
}

\setmainfont{PT Sans}
\setsansfont{PT Sans}
\setmonofont{Andale Mono}[Scale=MatchLowercase]
\renewcommand{\familydefault}{\sfdefault}
\setlength{\parindent}{0pt}
\setlength{\parskip}{0.56em}
\linespread{1.08}

\hypersetup{
  colorlinks=true,
  linkcolor=piaccent,
  urlcolor=piaccent,
  citecolor=piaccent,
  pdftitle={Pi Coding Agent},
  pdfauthor={Julius Darang}
}

% Headings: restrained, orange chapter markers, no oversized slide typography.
\titleformat{\section}
  {\Large\sffamily\bfseries\color{piaccent}}
  {}{0pt}{}
\titleformat{\subsection}
  {\large\sffamily\bfseries\color{pibody}}
  {}{0pt}{}
\titleformat{\subsubsection}
  {\normalsize\sffamily\bfseries\color{pimut}}
  {}{0pt}{}
\titlespacing*{\section}{0pt}{1.5em}{0.65em}
\titlespacing*{\subsection}{0pt}{1.05em}{0.35em}
\titlespacing*{\subsubsection}{0pt}{0.8em}{0.2em}

% Page furniture.
\pagestyle{fancy}
\fancyhf{}
\fancyhead[L]{\footnotesize\sffamily\color{pimut}PI CODING AGENT}
\fancyhead[R]{\footnotesize\sffamily\color{pimut}JULIUS DARANG}
\fancyfoot[C]{\footnotesize\sffamily\color{pimut}\thepage}
\renewcommand{\headrulewidth}{0.35pt}
\renewcommand{\footrulewidth}{0pt}
\makeatletter
\renewcommand{\maketitle}{%
  \begin{titlepage}
    \newgeometry{margin=0.9in}
    \pagecolor{piblack}\color{white}
    \vspace*{1.2in}
    {\sffamily\footnotesize\color{piaccent}\bfseries PI CODING AGENT\par}
    \vspace{0.45in}
    {\sffamily\fontsize{34}{40}\selectfont\bfseries\color{white}\@title\par}
    \vspace{0.28in}
    {\sffamily\Large\color{white!72}Build your coding workflow around a small terminal agent\par}
    \vfill
    {\sffamily\small\color{white!58}A practical guide to building a safe, project-aware terminal coding workflow.\par}
    \vspace{0.18in}
    {\sffamily\small\color{piaccent}\@author\par}
    \restoregeometry
  \end{titlepage}
  \pagecolor{white}\color{pibody}
  \pagestyle{fancy}
}
\makeatother

% Full-page chapter dividers for the four-part tutorial structure.
\newcommand{\partpage}[3]{%
  \clearpage
  \phantomsection
  \addcontentsline{toc}{section}{Part #1 -- #2}
  \pdfbookmark[1]{Part #1 -- #2}{part-#1}
  \thispagestyle{empty}
  \pagecolor{piblack}\color{white}
  \begin{center}
    \vspace*{1.35in}
    {\sffamily\footnotesize\color{piaccent}\bfseries PART #1\par}
    \vspace{0.42in}
    {\sffamily\fontsize{30}{36}\selectfont\bfseries\color{white}#2\par}
    \vspace{0.32in}
    {\sffamily\large\color{white!70}#3\par}
    \vfill
    {\sffamily\footnotesize\color{white!45}PI CODING AGENT · JULIUS DARANG\par}
  \end{center}
  \clearpage
  \pagecolor{white}\color{pibody}
  \pagestyle{fancy}
}

% Code blocks are readable on paper and retain the terminal character.
\definecolor{shadecolor}{HTML}{F3F0EB}
\usepackage{fvextra}
\DefineVerbatimEnvironment{Highlighting}{Verbatim}{fontsize=\small,breaklines=true,breakanywhere=true,commandchars=\\\{\}}
\setlength{\fboxsep}{8pt}

% Lists should breathe without becoming slide-like.
\setlist[itemize]{leftmargin=1.35em,itemsep=0.18em,topsep=0.2em}
\setlist[enumerate]{leftmargin=1.55em,itemsep=0.18em,topsep=0.2em}

% Make blockquotes read as practical callouts.
\renewenvironment{quote}{%
  \begin{leftbar}\small\color{pimut}
}{%
  \end{leftbar}
}
\renewcommand{\FrameCommand}{\hspace{0.4em}\textcolor{piaccent}{\vrule width 1.5pt}\hspace{0.8em}}

% Keep the table of contents compact and usable.
\renewcommand{\contentsname}{Contents}
\setcounter{tocdepth}{2}
\let\pioriginaltableofcontents\tableofcontents
\renewcommand{\tableofcontents}{\pioriginaltableofcontents\clearpage}
